及
\[
\mvfit{$\displaystyle\begin{aligned}
\begin{bmatrix}\hat{x}_{41}\\\hat{x}_{42}\end{bmatrix}(x_{43})
&=E\left(\begin{bmatrix}X_{41}X_{43}\\X_{42}X_{43}\end{bmatrix}\mid x_{43}=5;\boldsymbol{\mu},\boldsymbol{\Sigma}\right)
=\begin{bmatrix}\hat{x}_{41}\\\hat{x}_{42}\end{bmatrix}(x_{43})\\
&=\begin{bmatrix}6.4\\1.3\end{bmatrix}\times5=\begin{bmatrix}32.0\\6.5\end{bmatrix}
\end{aligned}$}
\]
于是, 我们预测出完全数据充分统计量为:
\[
\mvfit{$\displaystyle\begin{aligned}
\bar{\mathbf{T}}_1&=\begin{bmatrix}
\hat{x}_{11}+x_{21}+x_{31}+\hat{x}_{41}\\
x_{12}+x_{22}+x_{32}+\hat{x}_{42}\\
x_{13}+x_{23}+x_{33}+x_{43}
\end{bmatrix}
=\begin{bmatrix}
5.73+7+5+6.4\\
0+2+1+1.3\\
3+6+2+5
\end{bmatrix}
=\begin{bmatrix}24.13\\4.30\\16.00\end{bmatrix}
\end{aligned}$}
\]
\[
\mvfit{$\displaystyle\begin{aligned}
\bar{\mathbf{T}}_2&=\begin{bmatrix}
\hat{x}_{11}^2+x_{21}^2+x_{31}^2+\hat{x}_{41}^2 & \hat{x}_{11}x_{12}+x_{21}x_{22}+x_{31}x_{32}+\hat{x}_{41}\hat{x}_{42} & \hat{x}_{11}x_{13}+x_{21}x_{23}+x_{31}x_{33}+\hat{x}_{41}x_{43}\\
\hat{x}_{11}x_{12}+x_{21}x_{22}+x_{31}x_{32}+\hat{x}_{41}\hat{x}_{42} & x_{12}^2+x_{22}^2+x_{32}^2+\hat{x}_{42}^2 & x_{12}x_{13}+x_{22}x_{23}+x_{32}x_{33}+\hat{x}_{42}x_{43}\\
\hat{x}_{11}x_{13}+x_{21}x_{23}+x_{31}x_{33}+\hat{x}_{41}x_{43} & x_{12}x_{13}+x_{22}x_{23}+x_{32}x_{33}+\hat{x}_{42}x_{43} & x_{13}^2+x_{23}^2+x_{33}^2+x_{43}^2
\end{bmatrix}\\
&=\begin{bmatrix}
32.99+7^2+5^2+41.06 & 0+7\times2+5\times1+8.27 & 17.18+7\times6+5\times2+32\\
0+7\times2+5\times1+8.27 & 0+2^2+1^2+1.97 & 0\times3+2\times6+1\times2+6.5\\
17.18+7\times6+5\times2+32 & 0\times3+2\times6+1\times2+6.5 & 3^2+6^2+2^2+5^2
\end{bmatrix}\\
&=\begin{bmatrix}
148.05&27.27&101.18\\
27.27&6.97&20.50\\
101.18&20.50&74.00
\end{bmatrix}
\end{aligned}$}
\]
预测步骤到此结束.

在接下来的估计步骤中, 用\hyperref[eq:5-40]{式(5-40)}求出 $\boldsymbol{\mu}$ 与 $\boldsymbol{\Sigma}$ 的修正估计\footnote{最后形式的 $\boldsymbol{\Sigma}$ 的元素严格取两位小数.}:
\[
\mvfit{$\displaystyle\begin{aligned}
\hat{\boldsymbol{\mu}}&=\frac{1}{n}\bar{\mathbf{T}}_1=\frac{1}{4}
\begin{bmatrix}24.13\\4.30\\16.00\end{bmatrix}
=\begin{bmatrix}6.03\\1.08\\4.00\end{bmatrix}\\
\hat{\boldsymbol{\Sigma}}&=\frac{1}{n}\bar{\mathbf{T}}_2-\hat{\boldsymbol{\mu}}\hat{\boldsymbol{\mu}}'\\
&=\frac{1}{4}\begin{bmatrix}
148.05&27.27&101.18\\
27.27&6.97&20.50\\
101.18&20.50&74.00
\end{bmatrix}
-\begin{bmatrix}6.03\\1.08\\4.00\end{bmatrix}
\begin{bmatrix}6.03&1.08&4.00\end{bmatrix}\\
&=\begin{bmatrix}
0.61&0.33&1.17\\
0.33&0.59&0.83\\
1.17&0.83&2.50
\end{bmatrix}
\end{aligned}$}
\]
注意, 这里的 $\hat{\sigma}_{11}=0.61$ 与 $\hat{\sigma}_{22}=0.59$ 要比相应的初始估计大一些, 而初始估计是第一个变量和第二个变量中的缺损观测值用其余诸值的样本均值代替后得到的. 第三个方差的估计 $\hat{\sigma}_{33}$ 保持不变, 因为它不受缺损分量的影响.

预测步骤与估计步骤之间的迭代一直持续到 $\boldsymbol{\mu}$ 与 $\boldsymbol{\Sigma}$ 中的元素基本保持不变为止. 这类计算用计算机是很容易处理的. \hfill$\blacksquare$
